\section{Integer compositions}

\subsection{Compositions}

\begin{definition}
\label{def.fps.comps}
\lean{AlgebraicCombinatorics.Composition, AlgebraicCombinatorics.Composition.size, AlgebraicCombinatorics.Composition.len, AlgebraicCombinatorics.Composition.ofSize, AlgebraicCombinatorics.Composition.ofSizeIntoParts}
\leantarget
\leanok
\textbf{(a)} An \emph{(integer) composition} means a
(finite) tuple of positive integers. \medskip

\textbf{(b)} The \emph{size} of an integer composition $\alpha=\left(
\alpha_{1},\alpha_{2},\ldots,\alpha_{m}\right)  $ is defined to be the integer
$\alpha_{1}+\alpha_{2}+\cdots+\alpha_{m}$. It is written $\left\vert
\alpha\right\vert $. \medskip

\textbf{(c)} The \emph{length} of an integer composition $\alpha=\left(
\alpha_{1},\alpha_{2},\ldots,\alpha_{m}\right)  $ is defined to be the integer
$m$. It is written $\ell\left(  \alpha\right)  $. \medskip

\textbf{(d)} Let $n\in\mathbb{N}$. A \emph{composition of }$n$ means a
composition whose size is $n$. \medskip

\textbf{(e)} Let $n\in\mathbb{N}$ and $k\in\mathbb{N}$. A \emph{composition of
}$n$\emph{ into }$k$\emph{ parts} is a composition whose size is $n$ and whose
length is $k$.
\end{definition}

\begin{lemma}
\label{lem.fps.comps.size-nil}
\lean{AlgebraicCombinatorics.Composition.size_nil}
\leanhelper
\leanok
The size of the empty composition is $0$.
\end{lemma}

\begin{proof}
\leanok
By definition.
\end{proof}

\begin{lemma}
\label{lem.fps.comps.len-nil}
\lean{AlgebraicCombinatorics.Composition.len_nil}
\leanhelper
\leanok
The length of the empty composition is $0$.
\end{lemma}

\begin{proof}
\leanok
By definition.
\end{proof}

\begin{lemma}
\label{lem.fps.comps.size-cons}
\lean{AlgebraicCombinatorics.Composition.size_cons}
\leanhelper
\leanok
The size of a composition $(a) \cdot \alpha$ (prepending $a$ to $\alpha$)
is $a + |\alpha|$.
\end{lemma}

\begin{proof}
\leanok
By definition.
\end{proof}

\begin{lemma}
\label{lem.fps.comps.len-cons}
\lean{AlgebraicCombinatorics.Composition.len_cons}
\leanhelper
\leanok
The length of a composition $(a) \cdot \alpha$ (prepending $a$ to $\alpha$)
is $\ell(\alpha) + 1$.
\end{lemma}

\begin{proof}
\leanok
By definition of length.
\end{proof}

\begin{lemma}
\label{lem.fps.comps.size-ge-len}
\lean{AlgebraicCombinatorics.Composition.size_ge_len}
\leanhelper
\leanok
The size of a composition is at least its length: $\ell(\alpha) \leq |\alpha|$,
since each part is a positive integer (at least $1$).
\end{lemma}

\begin{proof}
\leanok
By induction on the composition. The base case (empty composition) is trivial.
For the inductive step, $|\alpha| = a + |\beta| \geq 1 + \ell(\beta) = \ell(\alpha)$
since $a \geq 1$.
\end{proof}

\begin{lemma}
\label{lem.fps.comps.empty-of-size-zero}
\lean{AlgebraicCombinatorics.Composition.empty_of_size_zero}
\leanhelper
\leanok
The empty list is the only composition of size $0$.
\end{lemma}

\begin{proof}
\leanok
If $\alpha = (a) \cdot \beta$, then $|\alpha| = a + |\beta| \geq 1 > 0$
since $a$ is a positive integer. Hence a composition of size $0$ must be empty.
\end{proof}

\begin{lemma}
\label{lem.fps.comps.equiv-mathlib}
\lean{AlgebraicCombinatorics.Composition.equivMathlib}
\leanhelper
\leanok
The definition of compositions of $n$ from Definition~\ref{def.fps.comps}
(as lists of positive integers with given sum) is equivalent to the definition
from Mathlib's composition type.
\end{lemma}

\begin{proof}
\leanok
Convert between lists of positive integers and lists of natural numbers
with positivity proofs.
\end{proof}

\begin{lemma}
\label{lem.fps.comps.equiv-mathlib-filtered}
\lean{AlgebraicCombinatorics.Composition.equivMathlibFiltered}
\leanhelper
\leanok
There is an equivalence between compositions of $n$ into $k$ parts
and the subtype of Mathlib compositions of $n$ having length $k$.
\end{lemma}

\begin{proof}
\leanok
Follows from the equivalence in Lemma~\ref{lem.fps.comps.equiv-mathlib},
restricting to compositions of the appropriate length.
\end{proof}

\subsection{Helpers for Theorem~\ref{thm.fps.comps.num-comps-n-k}}

\begin{lemma}
\label{lem.fps.comps.compositionToFinset}
\lean{AlgebraicCombinatorics.MathlibComposition.compositionToFinset}
\leanhelper
\leanok
There is an equivalence between Mathlib's compositions of $n$ and
subsets of $\{1, 2, \ldots, n-1\}$,
obtained by composing two Mathlib equivalences.
\end{lemma}

\begin{proof}
\leanok
This is the composition of the Mathlib equivalences
between compositions of $n$, composition-as-set representations,
and finite subsets of $\{1, 2, \ldots, n-1\}$.
\end{proof}

\begin{lemma}
\label{lem.fps.comps.compositionAsSetEquiv-card}
\lean{AlgebraicCombinatorics.MathlibComposition.compositionAsSetEquiv_card_eq_length_sub_one}
\leanhelper
\leanok
For $n > 0$ and a composition-as-set $c$ of $n$, the cardinality of the
associated finset of interior boundary points $S(c)$ satisfies
$|S(c)| + 1 = \ell(c)$.
This holds because the equivalence extracts the interior boundary points
(those other than $0$ and $n$), and a composition of length $k$ has exactly
$k - 1$ such points.
\end{lemma}

\begin{proof}
The boundary set of $c$ has $\ell(c) + 1$ elements (including $0$ and $n$).
The equivalence maps to the interior points, obtained by removing $0$ and $n$.
The resulting finset has $\ell(c) + 1 - 2 = \ell(c) - 1$ elements.
\end{proof}

\begin{lemma}
\label{lem.fps.comps.compositionToFinset-card}
\lean{AlgebraicCombinatorics.MathlibComposition.compositionToFinset_card_add_one}
\leanhelper
\leanok
For $n > 0$ and a composition $c$ of $n$, the finset $S(c)$ corresponding to $c$
under the composition-to-finset equivalence satisfies
$|S(c)| + 1 = \ell(c)$.
\end{lemma}

\begin{proof}
\leanok
Follows directly from Lemma~\ref{lem.fps.comps.compositionAsSetEquiv-card}
applied to the composition-as-set representation of $c$.
\end{proof}

\begin{lemma}
\label{lem.fps.comps.card-compositions-of-length}
\lean{AlgebraicCombinatorics.MathlibComposition.card_compositions_of_length}
\leanhelper
\leanok
For $n > 0$ and $k > 0$, the number of compositions of $n$ into $k$ parts
(using Mathlib's composition type) equals $\binom{n-1}{k-1}$.
The proof establishes a bijection between compositions of length $k$
and $(k-1)$-element subsets of $\{1, 2, \ldots, n-1\}$ via
Lemma~\ref{lem.fps.comps.compositionToFinset-card}, then uses
the formula for counting subsets of a given size.
\end{lemma}

\begin{proof}
Compositions of length $k$ correspond (via the composition-to-finset equivalence) to
$(k-1)$-element subsets of $\{1, 2, \ldots, n-1\}$.
These are counted by $\binom{n-1}{k-1}$.
\end{proof}

\begin{theorem}
\label{thm.fps.comps.num-comps-n-k}
\lean{AlgebraicCombinatorics.Composition.card_ofSizeIntoParts_pos, AlgebraicCombinatorics.Composition.card_ofSizeIntoParts_k_zero, AlgebraicCombinatorics.Composition.card_ofSizeIntoParts_zero}
\leantarget
\leanok
Let $n,k\in\mathbb{N}$. Then, the \# of
compositions of $n$ into $k$ parts is%
\[
\dbinom{n-1}{n-k}=%
\begin{cases}
\dbinom{n-1}{k-1}, & \text{if }n>0;\\
\delta_{k,0}, & \text{if }n=0.
\end{cases}
\]
\end{theorem}

\begin{proof}
Fix $k\in\mathbb{N}$, but don't fix $n$. Let
\begin{equation}
a_{n,k}=\left(  \text{\# of compositions of }n\text{ into }k\text{
parts}\right)  . \label{pf.thm.fps.comps.n-into-k-parts.1st.ank=}%
\end{equation}
We want to find $a_{n,k}$. We define the generating function%
\begin{equation}
A_{k}:=\sum_{n\in\mathbb{N}}a_{n,k}x^{n}=\left(  a_{0,k},a_{1,k}%
,a_{2,k},\ldots\right)  \in\mathbb{Q}\left[  \left[  x\right]  \right]  .
\label{pf.thm.fps.comps.n-into-k-parts.1st.Ak=}%
\end{equation}

Let us write $\mathbb{P}$ for the set $\left\{  1,2,3,\ldots\right\}  $. Then,
a composition of $n$ into $k$ parts is nothing but a $k$-tuple $\left(
\alpha_{1},\alpha_{2},\ldots,\alpha_{k}\right)  \in\mathbb{P}^{k}$ satisfying
$\alpha_{1}+\alpha_{2}+\cdots+\alpha_{k}=n$. Hence,
(\ref{pf.thm.fps.comps.n-into-k-parts.1st.ank=}) can be rewritten as%
\begin{align}
a_{n,k}  &  =\left(  \text{\# of all }k\text{-tuples }\left(  \alpha
_{1},\alpha_{2},\ldots,\alpha_{k}\right)  \in\mathbb{P}^{k}\text{ satisfying
}\alpha_{1}+\alpha_{2}+\cdots+\alpha_{k}=n\right) \nonumber\\
&  =\sum\limits_{\substack{\left(  \alpha_{1},\alpha_{2},\ldots,\alpha
_{k}\right)  \in\mathbb{P}^{k};\\\alpha_{1}+\alpha_{2}+\cdots+\alpha_{k}=n}}1.
\label{pf.thm.fps.comps.n-into-k-parts.1st.ank=2}%
\end{align}
Thus, we can rewrite the equality $A_{k}=\sum_{n\in\mathbb{N}}a_{n,k}x^{n}$
as
\begin{align*}
A_{k}  &  =\sum_{n\in\mathbb{N}}\left(  \sum\limits_{\substack{\left(
\alpha_{1},\alpha_{2},\ldots,\alpha_{k}\right)  \in\mathbb{P}^{k};\\\alpha
_{1}+\alpha_{2}+\cdots+\alpha_{k}=n}}1\right)  x^{n}=\sum_{n\in\mathbb{N}%
}\ \ \sum\limits_{\substack{\left(  \alpha_{1},\alpha_{2},\ldots,\alpha
_{k}\right)  \in\mathbb{P}^{k};\\\alpha_{1}+\alpha_{2}+\cdots+\alpha_{k}%
=n}}\ \ \underbrace{x^{n}}_{\substack{=x^{\alpha_{1}+\alpha_{2}+\cdots
+\alpha_{k}}\\\text{(since }\alpha_{1}+\alpha_{2}+\cdots+\alpha_{k}=n\text{)}%
}}\\
&  =\underbrace{\sum_{n\in\mathbb{N}}\ \ \sum\limits_{\substack{\left(
\alpha_{1},\alpha_{2},\ldots,\alpha_{k}\right)  \in\mathbb{P}^{k};\\\alpha
_{1}+\alpha_{2}+\cdots+\alpha_{k}=n}}}_{=\sum\limits_{\left(  \alpha
_{1},\alpha_{2},\ldots,\alpha_{k}\right)  \in\mathbb{P}^{k}}}%
\ \ \underbrace{x^{\alpha_{1}+\alpha_{2}+\cdots+\alpha_{k}}}_{=x^{\alpha_{1}%
}x^{\alpha_{2}}\cdots x^{\alpha_{k}}}=\sum\limits_{\left(  \alpha_{1}%
,\alpha_{2},\ldots,\alpha_{k}\right)  \in\mathbb{P}^{k}}x^{\alpha_{1}%
}x^{\alpha_{2}}\cdots x^{\alpha_{k}}\\
&  =\left(  \sum_{\alpha_{1}\in\mathbb{P}}x^{\alpha_{1}}\right)  \left(
\sum_{\alpha_{2}\in\mathbb{P}}x^{\alpha_{2}}\right)  \cdots\left(
\sum_{\alpha_{k}\in\mathbb{P}}x^{\alpha_{k}}\right) \\
&  =\left(  \sum_{n\in\mathbb{P}}x^{n}\right)  ^{k}%
\end{align*}
(here, we have renamed all the $k$ summation indices as $n$, and realized that
all $k$ sums are identical). Since%
\[
\sum_{n\in\mathbb{P}}x^{n}=x^{1}+x^{2}+x^{3}+\cdots=x\underbrace{\left(
1+x+x^{2}+\cdots\right)  }_{=\dfrac{1}{1-x}}=x\cdot\dfrac{1}{1-x}=\dfrac
{x}{1-x},
\]
this can be rewritten further as%
\begin{equation}
A_{k}=\left(  \dfrac{x}{1-x}\right)  ^{k}=x^{k}\left(  1-x\right)  ^{-k}.
\label{pf.thm.fps.comps.num-comps-n-k.Ak=}%
\end{equation}

In order to simplify this, we need to expand $\left(  1-x\right)  ^{-k}$. This
is routine by now: Theorem \ref{thm.fps.newton-binom} (applied to $-k$ instead
of $n$) yields%
\[
\left(  1+x\right)  ^{-k}=\sum_{j\in\mathbb{N}}\dbinom{-k}{j}x^{j}.
\]
Substituting $-x$ for $x$ on both sides of this equality (i.e., applying the
map $f\mapsto f\circ\left(  -x\right)  $), we obtain%
\begin{align}
\left(  1-x\right)  ^{-k}  &  =\sum_{j\in\mathbb{N}}\underbrace{\dbinom{-k}%
{j}}_{=\left(  -1\right)  ^{j}\dbinom{j+k-1}{j}}\underbrace{\left(  -x\right)  ^{j}}_{=\left(
-1\right)  ^{j}x^{j}}=\sum_{j\in\mathbb{N}}\left(  -1\right)  ^{j}%
\dbinom{j+k-1}{j}\cdot\left(  -1\right)  ^{j}x^{j}\nonumber\\
&  =\sum_{j\in\mathbb{N}}\underbrace{\left(  \left(  -1\right)  ^{j}\right)
^{2}}_{=1}\dbinom{j+k-1}{j}x^{j}=\sum_{j\in\mathbb{N}}\dbinom{j+k-1}{j}%
x^{j}\label{pf.thm.fps.comps.num-comps-n-k.1-x-k}\\
&  =\sum\limits_{n\geq k}\dbinom{n-1}{n-k}x^{n-k}\nonumber
\end{align}
(where we used Theorem \ref{thm.binom.upneg-n}, applied to $k$ and $j$ instead of $n$ and $k$, to simplify $\dbinom{-k}{j}$;
and we have substituted $n-k$ for $j$ in the sum). Hence, our above
computation of $A_{k}$ can be completed as follows:
\begin{align*}
A_{k}  &  =x^{k}\underbrace{\left(  1-x\right)  ^{-k}}_{=\sum\limits_{n\geq
k}\dbinom{n-1}{n-k}x^{n-k}}=x^{k}\sum\limits_{n\geq k}\dbinom{n-1}{n-k}%
x^{n-k}\\
&  =\sum\limits_{n\geq k}\dbinom{n-1}{n-k}\underbrace{x^{k}x^{n-k}}_{=x^{n}%
}=\sum\limits_{n\geq k}\dbinom{n-1}{n-k}x^{n}=\sum\limits_{n\in\mathbb{N}%
}\dbinom{n-1}{n-k}x^{n}%
\end{align*}
(here, we have extended the range of the summation from all $n\geq k$ to all
$n\in\mathbb{N}$; this did not change the sum, since all the newly introduced
addends with $n<k$ are $0$). Comparing coefficients, we thus obtain%
\begin{equation}
a_{n,k}=\dbinom{n-1}{n-k}\ \ \ \ \ \ \ \ \ \ \text{for each }n\in\mathbb{N}.
\label{pf.thm.fps.comps.n-into-k-parts.1st.rhs1}%
\end{equation}

If $n>0$, then we can rewrite the right hand side of this equality as
$\dbinom{n-1}{k-1}$ (using Theorem \ref{thm.binom.sym}). However, if $n=0$,
then this right hand side equals $\delta_{k,0}$ instead (where we are using
Definition \ref{def.kron-delta}). Thus, we can rewrite
(\ref{pf.thm.fps.comps.n-into-k-parts.1st.rhs1}) as%
\begin{equation}
a_{n,k}=%
\begin{cases}
\dbinom{n-1}{k-1}, & \text{if }n>0;\\
\delta_{k,0}, & \text{if }n=0
\end{cases}
\ \ \ \ \ \ \ \ \ \ \text{for each }n\in\mathbb{N}.
\label{pf.thm.fps.comps.n-into-k-parts.1st.rhs2}%
\end{equation}

In the Lean formalization, this proceeds by:
\begin{enumerate}
\item Using the equivalence between our compositions and Mathlib's
composition type (Lemma~\ref{lem.fps.comps.equiv-mathlib-filtered}).
\item Using Mathlib's equivalence between compositions of $n$ and subsets of
$\{1, 2, \ldots, n-1\}$.
\item Showing that compositions of length $k$ correspond to subsets of size $k-1$.
\item Computing $|\{S \subseteq \{1, \ldots, n-1\} \mid |S| = k-1\}| = \binom{n-1}{k-1}$
using the formula for counting subsets of a given size.
\end{enumerate}

The edge cases are:
\begin{itemize}
\item For $n > 0$ and $k = 0$: There are no compositions (since parts are positive,
and a tuple of zero positive integers sums to $0$, not $n$). This gives $0 = \binom{n-1}{-1} = 0$.
\item For $n = 0$: The empty tuple is the only composition of $0$ (by
Lemma~\ref{lem.fps.comps.empty-of-size-zero}), and it has length $0$.
So the count is $\delta_{k,0}$.
\end{itemize}
\end{proof}

\begin{theorem}
\label{thm.fps.comps.num-comps-n}
\lean{AlgebraicCombinatorics.Composition.card_ofSize, AlgebraicCombinatorics.Composition.card_ofSize_pos}
\leantarget
\leanok
Let $n\in\mathbb{N}$. Then, the \# of
compositions of $n$ is%
\[%
\begin{cases}
2^{n-1}, & \text{if }n>0;\\
1, & \text{if }n=0.
\end{cases}
\]
\end{theorem}

\begin{proof}
If $n=0$, then
the \# of compositions of $n$ is $1$ (since the empty tuple $\left(
{}\right)  $ is the only composition of $0$). Thus, for the rest of this
proof, we WLOG assume that $n\neq0$. Hence, we must prove that the \# of
compositions of $n$ is $2^{n-1}$.

If $\left(  n_{1},n_{2},\ldots,n_{m}\right)  $ is a composition of $n$, then
$m\in\left\{  1,2,\ldots,n\right\}  $. In other words, any composition
of $n$ is a composition of $n$ into $k$ parts for some $k\in\left\{
1,2,\ldots,n\right\}  $. Hence,%
\begin{align*}
&  \left(  \text{\# of compositions of }n\right) \\
&  =\sum_{k\in\left\{  1,2,\ldots,n\right\}  }\underbrace{\left(  \text{\# of
compositions of }n\text{ into }k\text{ parts}\right)  }_{=\dbinom
{n-1}{n-k}}\\
&  =\sum_{k\in\left\{  1,2,\ldots,n\right\}  }\dbinom{n-1}{n-k}=\sum_{k=1}%
^{n}\dbinom{n-1}{n-k}=\sum_{k=0}^{n-1}\dbinom{n-1}{k}%
\end{align*}
(where the second equality uses Theorem \ref{thm.fps.comps.num-comps-n-k};
and we have substituted $k$ for $n-k$ in the sum). Comparing this with%
\begin{align*}
2^{n-1}  &  =\left(  1+1\right)  ^{n-1}=\sum_{k=0}^{n-1}\dbinom{n-1}%
{k}\underbrace{1^{k}1^{n-1-k}}_{=1}\ \ \ \ \ \ \ \ \ \ \left(  \text{by the
binomial theorem}\right) \\
&  =\sum_{k=0}^{n-1}\dbinom{n-1}{k},
\end{align*}
we obtain $\left(  \text{\# of compositions of }n\right)  =2^{n-1}$.

In the Lean formalization, this is proved by establishing an equivalence with
Mathlib's composition type (Lemma~\ref{lem.fps.comps.equiv-mathlib})
and using Mathlib's theorem on the cardinality of compositions.
\end{proof}

\subsection{Weak compositions}

\begin{definition}
\label{def.fps.wcomps}
\lean{AlgebraicCombinatorics.WeakComposition, AlgebraicCombinatorics.WeakComposition.size, AlgebraicCombinatorics.WeakComposition.len, AlgebraicCombinatorics.WeakComposition.ofSize, AlgebraicCombinatorics.WeakComposition.ofSizeIntoParts}
\leantarget
\leanok
\textbf{(a)} An \emph{(integer) weak composition} means
a (finite) tuple of nonnegative integers. \medskip

\textbf{(b)} The \emph{size} of a weak composition $\alpha=\left(  \alpha
_{1},\alpha_{2},\ldots,\alpha_{m}\right)  $ is defined to be the integer
$\alpha_{1}+\alpha_{2}+\cdots+\alpha_{m}$. It is written $\left\vert
\alpha\right\vert $. \medskip

\textbf{(c)} The \emph{length} of a weak composition $\alpha=\left(
\alpha_{1},\alpha_{2},\ldots,\alpha_{m}\right)  $ is defined to be the integer
$m$. It is written $\ell\left(  \alpha\right)  $. \medskip

\textbf{(d)} Let $n\in\mathbb{N}$. A \emph{weak composition of }$n$ means a
weak composition whose size is $n$. \medskip

\textbf{(e)} Let $n\in\mathbb{N}$ and $k\in\mathbb{N}$. A \emph{weak
composition of }$n$\emph{ into }$k$\emph{ parts} is a weak composition whose
size is $n$ and whose length is $k$.
\end{definition}

\begin{lemma}
\label{lem.fps.wcomps.equiv-fun}
\lean{AlgebraicCombinatorics.WeakComposition.ofSizeIntoParts_equiv}
\leanhelper
\leanok
The list-based representation of weak compositions of $n$ into $k$ parts
is equivalent to the function-based representation
$\{f : \{0, 1, \ldots, k-1\} \to \mathbb{N} \mid \sum_{i} f(i) = n\}$.
\end{lemma}

\begin{proof}
\leanok
The forward direction converts a list $[\alpha_0, \alpha_1, \ldots, \alpha_{k-1}]$
to the function $i \mapsto \alpha_i$.
The backward direction converts a function $f$
to the list $[f(0), f(1), \ldots, f(k-1)]$.
\end{proof}

\begin{lemma}
\label{lem.fps.wcomps.equiv-sym}
\lean{AlgebraicCombinatorics.WeakComposition.ofSizeIntoParts'.equivSym}
\leanhelper
\leanok
The function-based weak compositions of $n$ into $k$ parts are equivalent
to multisets of size $n$ from a $k$-element alphabet.
\end{lemma}

\begin{proof}
\leanok
The bijection sends $f : \{0, \ldots, k-1\} \to \mathbb{N}$ with $\sum_i f(i) = n$
to the multiset containing $f(i)$ copies of $i$ for each $i$, and the inverse sends a
multiset $m$ to the function $i \mapsto \text{count}(i, m)$.
\end{proof}

\begin{lemma}
\label{lem.fps.wcomps.to-composition}
\lean{AlgebraicCombinatorics.WeakComposition.toComposition}
\leanhelper
\leanok
There is a bijection between weak compositions of $n$ into $k$ parts
and compositions of $n+k$ into $k$ parts, defined by adding $1$ to each entry.
\end{lemma}

\begin{proof}
\leanok
If $\left(  \alpha_{1},\alpha_{2},\ldots,\alpha_{k}\right)  $ is a weak
composition of $n$ into $k$ parts, then $\left(  \alpha_{1}%
+1,\alpha_{2}+1,\ldots,\alpha_{k}+1\right)  $ is a composition of $n+k$ into
$k$ parts (since adding $1$ to each nonnegative entry yields a positive integer,
and the sum increases by $k$).
The inverse map subtracts $1$ from each entry.
\end{proof}

\begin{lemma}
\label{lem.fps.wcomps.card-eq}
\lean{AlgebraicCombinatorics.WeakComposition.ofSizeIntoParts'.card_eq}
\leanhelper
\leanok
The number of function-based weak compositions of $n$ into $k$ parts
equals $\binom{n+k-1}{n}$. The proof uses the equivalence with
multisets of size $n$ from a $k$-element alphabet (Lemma~\ref{lem.fps.wcomps.equiv-sym})
and the formula for counting such multisets.
\end{lemma}

\begin{proof}
\leanok
By the equivalence with multisets of size $n$ from a $k$-element alphabet and the formula
$\binom{n + k - 1}{n}$ for counting such multisets.
\end{proof}

\begin{theorem}
\label{thm.fps.comps.num-wcomps-n-k}
\lean{AlgebraicCombinatorics.WeakComposition.card_ofSizeIntoParts, AlgebraicCombinatorics.WeakComposition.card_ofSizeIntoParts_pos, AlgebraicCombinatorics.WeakComposition.card_ofSizeIntoParts_zero}
\leantarget
\leanok
Let $n,k\in\mathbb{N}$. Then, the \# of
weak compositions of $n$ into $k$ parts is
\[
\dbinom{n+k-1}{n}=%
\begin{cases}
\dbinom{n+k-1}{k-1}, & \text{if }k>0;\\
\delta_{n,0}, & \text{if }k=0.
\end{cases}
\]
\end{theorem}

\begin{proof}
Adding $1$
to each entry of a weak composition of $n$ into $k$ parts gives
a composition of $n+k$ into $k$ parts (Lemma~\ref{lem.fps.wcomps.to-composition}).
This map is a bijection (its inverse subtracts $1$ from each entry).
Thus, by the bijection principle, we
have%
\begin{align*}
&  \left\vert \left\{  \text{weak compositions of }n\text{ into }k\text{
parts}\right\}  \right\vert \\
&  =\left\vert \left\{  \text{compositions of }n+k\text{ into }k\text{
parts}\right\}  \right\vert \\
&  =\left(  \text{\# of compositions of }n+k\text{ into }k\text{ parts}\right)
\\
&  =\dbinom{n+k-1}{n+k-k}
=\dbinom{n+k-1}{n}.
\end{align*}
(Here, the last step follows by Theorem \ref{thm.fps.comps.num-comps-n-k},
applied to $n+k$ instead of $n$.)
It remains to prove that this equals $%
\begin{cases}
\dbinom{n+k-1}{k-1}, & \text{if }k>0;\\
\delta_{n,0}, & \text{if }k=0
\end{cases}
$ as well. If $k=0$, then it is clear
by inspection; otherwise it follows from Theorem \ref{thm.binom.sym}.

In the Lean formalization, the primary proof uses the ``stars and bars'' technique:
weak compositions of $n$ into $k$ parts are equivalent to multisets of size $n$ from
a $k$-element alphabet, via
Lemmas~\ref{lem.fps.wcomps.equiv-fun} and~\ref{lem.fps.wcomps.equiv-sym}.
These are counted by $\binom{n+k-1}{n}$ using the formula for counting multisets.
\end{proof}

\subsection{Weak compositions with entries from $\left\{  0,1,\ldots
,p-1\right\}  $}

\begin{lemma}
\label{lem.fps.comps.geom-series-mul}
\lean{AlgebraicCombinatorics.BoundedWeakComposition.geom_series_mul_one_sub}
\leanhelper
\leanok
The geometric series identity: $\left(\sum_{i=0}^{p-1} x^i\right)(1 - x) = 1 - x^p$.
\end{lemma}

\begin{proof}
By induction on $p$ and ring arithmetic.
\end{proof}

\begin{lemma}
\label{lem.fps.comps.geom-series-eq-mul-inv}
\lean{AlgebraicCombinatorics.BoundedWeakComposition.geom_series_eq_mul_inv}
\leanhelper
\leanok
The geometric series equals $(1 - x^p) \cdot (1-x)^{-1}$:
\[
\sum_{i=0}^{p-1} x^i = (1 - x^p) \cdot (1-x)^{-1}.
\]
\end{lemma}

\begin{proof}
\leanok
Multiply both sides by $(1 - x)$ and use Lemma~\ref{lem.fps.comps.geom-series-mul}.
\end{proof}

\begin{lemma}
\label{lem.fps.comps.invOneSubPow-pow}
\lean{AlgebraicCombinatorics.BoundedWeakComposition.invOneSubPow_pow}
\leanhelper
\leanok
For all $k \in \mathbb{N}$,
$\left((1 - x)^{-1}\right)^k = (1 - x)^{-k}$.
That is, the $k$-th power of the formal inverse of $(1 - x)$ equals
the formal inverse of $(1 - x)^k$.
\end{lemma}

\begin{proof}
\leanok
By induction on $k$, using the identity $(1-x)^{-(a+b)} = (1-x)^{-a} \cdot (1-x)^{-b}$
from Mathlib.
\end{proof}

\begin{lemma}
\label{lem.fps.comps.coeff-invOneSubPow}
\lean{AlgebraicCombinatorics.BoundedWeakComposition.coeff_invOneSubPow}
\leanhelper
\leanok
For $k > 0$, the coefficient of $x^n$ in $(1-x)^{-k}$ is
$\binom{k - 1 + n}{k - 1}$.
\end{lemma}

\begin{proof}
\leanok
This is a standard result from Mathlib's power series library.
\end{proof}

\begin{lemma}
\label{lem.fps.comps.genFun-eq}
\lean{AlgebraicCombinatorics.BoundedWeakComposition.genFun_eq}
\leanhelper
\leanok
The generating function for bounded weak compositions satisfies
$W_{k,p} = (1-x^p)^k \cdot (1-x)^{-k}$.
\end{lemma}

\begin{proof}
\leanok
Use Lemma~\ref{lem.fps.comps.geom-series-eq-mul-inv} and raise both sides to the $k$-th power,
using Lemma~\ref{lem.fps.comps.invOneSubPow-pow} to simplify
$\left((1-x)^{-1}\right)^k = (1-x)^{-k}$.
\end{proof}

\begin{lemma}
\label{lem.fps.comps.coeff-genFun}
\lean{AlgebraicCombinatorics.BoundedWeakComposition.coeff_genFun}
\leanhelper
\leanok
The coefficient of $x^n$ in the generating function $W_{k,p}$ equals the number of
bounded weak compositions of $n$ into $k$ parts with entries in $\{0, 1, \ldots, p-1\}$.
\end{lemma}

\begin{proof}
\leanok
Use the product rule for power series coefficients
and show that the product of indicator functions picks out exactly the bounded
weak compositions summing to $n$.
\end{proof}

\begin{lemma}
\label{lem.fps.comps.coeff-genFun-formula}
\lean{AlgebraicCombinatorics.BoundedWeakComposition.coeff_genFun_formula}
\leanhelper
\leanok
The coefficient of $x^n$ in $W_{k,p}$ equals the inclusion-exclusion formula:
\[
[x^n] W_{k,p} = \sum_{\substack{j \in \{0,\ldots,k\} \\ pj \leq n}}
(-1)^j \binom{k}{j} \binom{n - pj + k - 1}{n - pj}.
\]
\end{lemma}

\begin{proof}
Expand the product $(1-x^p)^k \cdot (1-x)^{-k}$ using the binomial theorem for
$(1-x^p)^k$ and the expansion of $(1-x)^{-k}$ via
Lemma~\ref{lem.fps.comps.coeff-invOneSubPow},
then extract the coefficient of $x^n$ via the convolution formula.
\end{proof}

\begin{theorem}
\label{thm.fps.comps.num-wpcomps-n-k}
\lean{AlgebraicCombinatorics.BoundedWeakComposition.card}
\leantarget
\leanok
Let $n,k,p\in\mathbb{N}$. Then, the \# of
$k$-tuples $\left(  \alpha_{1},\alpha_{2},\ldots,\alpha_{k}\right)
\in\left\{  0,1,\ldots,p-1\right\}  ^{k}$ satisfying $\alpha_{1}+\alpha
_{2}+\cdots+\alpha_{k}=n$ is%
\[
\sum_{j=0}^{k}\left(  -1\right)  ^{j}\dbinom{k}{j}\dbinom{n-pj+k-1}{n-pj}.
\]
\end{theorem}

\begin{proof}
We fix three nonnegative integers $n$, $k$ and $p$. We are
looking for the \# of $k$-tuples $\left(  \alpha_{1},\alpha_{2},\ldots
,\alpha_{k}\right)  \in\left\{  0,1,\ldots,p-1\right\}  ^{k}$ satisfying
$\alpha_{1}+\alpha_{2}+\cdots+\alpha_{k}=n$. In other words, we are looking
for the \# of all weak compositions of $n$ into $k$ parts with the property
that each entry is $<p$. Let us denote this \# by $w_{n,k,p}$.

We invoke a generating function. Define the FPS%
\[
W_{k,p}:=\sum_{n\in\mathbb{N}}w_{n,k,p}x^{n}\in\mathbb{Q}\left[  \left[
x\right]  \right]  .
\]
For each $n\in\mathbb{N}$, we have%
\begin{align*}
w_{n,k,p}
&  =\sum\limits_{\substack{\left(  \alpha_{1},\alpha_{2},\ldots,\alpha
_{k}\right)  \in\left\{  0,1,\ldots,p-1\right\}  ^{k};\\\alpha_{1}+\alpha
_{2}+\cdots+\alpha_{k}=n}}1.
\end{align*}
Thus, we can rewrite $W_{k,p}=\sum_{n\in\mathbb{N}}w_{n,k,p}x^{n}$ as%
\begin{align*}
W_{k,p}
&  =\sum\limits_{\left(  \alpha_{1},\alpha_{2},\ldots,\alpha_{k}\right)
\in\left\{  0,1,\ldots,p-1\right\}  ^{k}}x^{\alpha_{1}}x^{\alpha_{2}}\cdots
x^{\alpha_{k}}\\
&  =\left(  \sum_{n\in\left\{  0,1,\ldots,p-1\right\}  }x^{n}\right)  ^{k}%
\end{align*}
(by the product rule). Since%
\[
\sum_{n\in\left\{  0,1,\ldots,p-1\right\}  }x^{n}=x^{0}+x^{1}+\cdots
+x^{p-1}=\dfrac{1-x^{p}}{1-x},
\]
this can be rewritten further as%
\begin{equation}
W_{k,p}=\left(  \dfrac{1-x^{p}}{1-x}\right)  ^{k}=\left(  1-x^{p}\right)
^{k}\left(  1-x\right)  ^{-k}. \label{pf.thm.fps.comps.num-wpcomps-n-k.3}%
\end{equation}
In order to expand the right hand side, let us expand $\left(  1-x^{p}\right)
^{k}$ and $\left(  1-x\right)  ^{-k}$ separately.

The binomial theorem yields%
\[
\left(  1-x^{p}\right)  ^{k}=\sum_{j\in\mathbb{N}}\left(  -1\right)
^{j}\dbinom{k}{j}x^{pj}.
\]
Multiplying this by
(\ref{pf.thm.fps.comps.num-comps-n-k.1-x-k}), we obtain%
\begin{align*}
&  \left(  1-x^{p}\right)  ^{k}\left(  1-x\right)  ^{-k}\\
&  =\sum_{n\in\mathbb{N}}\left(  \sum_{\substack{\left(  i,j\right)
\in\mathbb{N}\times\mathbb{N};\\pj+i=n}}\left(  -1\right)  ^{j}\dbinom{k}%
{j}\dbinom{i+k-1}{i}\right)  x^{n}.
\end{align*}
Thus, (\ref{pf.thm.fps.comps.num-wpcomps-n-k.3}) becomes
\[
W_{k,p}=\sum
_{n\in\mathbb{N}}\left(  \sum_{\substack{\left(  i,j\right)  \in
\mathbb{N}\times\mathbb{N};\\pj+i=n}}\left(  -1\right)  ^{j}\dbinom{k}%
{j}\dbinom{i+k-1}{i}\right)  x^{n}.
\]
Comparing coefficients, we find that each $n\in\mathbb{N}$
satisfies%
\begin{align*}
w_{n,k,p}
&  =\sum_{j=0}^{k}\left(  -1\right)  ^{j}\dbinom{k}{j}\dbinom{n-pj+k-1}{n-pj}.
\end{align*}
\end{proof}

\begin{lemma}
\label{lem.fps.comps.card-bwc-2}
\lean{AlgebraicCombinatorics.card_boundedWeakComposition_2}
\leanhelper
\leanok
Let $n,k\in\mathbb{N}$. The number of binary $k$-strings with exactly $n$
ones is $\binom{k}{n}$.
\end{lemma}

\begin{proof}
Establish a bijection between binary $k$-strings with $n$ ones and $n$-element
subsets of $\{1, 2, \ldots, k\}$ (a binary string corresponds to the set of positions
where a $1$ appears). Then use the formula for counting subsets of a given size.
\end{proof}

\begin{proposition}
\label{prop.fps.comps.num-w2comps-n-k-id}
\lean{AlgebraicCombinatorics.binom_sum_identity}
\leantarget
\leanok
Let $n,k\in\mathbb{N}$. Then,%
\[
\dbinom{k}{n}=\sum_{j=0}^{k}\left(  -1\right)  ^{j}\dbinom{k}{j}%
\dbinom{n-2j+k-1}{n-2j}.
\]
\end{proposition}

\begin{proof}
The $k$-tuples $\left(  \alpha_{1},\alpha
_{2},\ldots,\alpha_{k}\right)  \in\left\{  0,1\right\}  ^{k}$ are just the
\emph{binary }$k$\emph{-strings}, i.e., the
$k$-tuples formed of $0$s and $1$s. Imposing the condition $\alpha_{1}%
+\alpha_{2}+\cdots+\alpha_{k}=n$ on such a $k$-tuple is tantamount to
requiring that it contain exactly $n$ many $1$s. Therefore, the \# of all
$k$-tuples $\left(  \alpha_{1},\alpha_{2},\ldots,\alpha_{k}\right)
\in\left\{  0,1\right\}  ^{k}$ satisfying $\alpha_{1}+\alpha_{2}+\cdots
+\alpha_{k}=n$ is $\dbinom{k}{n}$, since we have to choose which $n$ of its
$k$ positions will be occupied by $1$s
(Lemma~\ref{lem.fps.comps.card-bwc-2}).
However, Theorem
\ref{thm.fps.comps.num-wpcomps-n-k} (applied to $p=2$) yields that this \#
equals $\sum_{j=0}^{k}\left(  -1\right)  ^{j}\dbinom{k}{j}\dbinom
{n-2j+k-1}{n-2j}$. Comparing these two results, we obtain the identity.
\end{proof}
